\chapter{The problem}\label{ch:problem}

\section{Casualty reporting assumes a known position}

Every casualty report format in use starts with a location. The NATO 9-line MEDEVAC request
puts it on line 1. Tactical apps and battle management systems take it from the device's GPS.
The challenge brief we answered (Voice-to-Map, sponsor Lysk) asks for spoken position updates
turned into map updates, and suggests OpenStreetMap geocoding plus an LLM for intent
extraction. Every team we saw built that pipeline: speech, then text, then a place name, then a
pin.

All of it fails in the same situation: the person reporting doesn't know where he is.

\section{GPS is jammed, and not only at the front}

Satellite navigation is now routinely denied or falsified over large areas, not only on the
front line. Governments around the Baltic reported to ICAO that about 123,000 flights were
affected by GPS interference between January and April 2025, with authorities in Estonia saying
85 \% of flights there were affected \cite{gps123k,trtgps}. Poland recorded 2,732 jamming and
spoofing cases in January 2025 alone \cite{trtgps}. A receiver under spoofing is worse than a
receiver with no fix: it shows a confident, wrong position.

On the ground, the soldier near the front knows roughly which sector he's in and can see what
is around him. He usually can't produce a grid reference under fire, at night or while treating
a wounded friend.

\section{Drones have removed the golden hour}

In Ukraine, persistent drone surveillance has changed how casualties are moved. A Ukrainian
officer at the NATO-Ukraine training centre said the golden hour ``does not exist'' any more, and
that only a ``golden day'' remains \cite{goldenday}. Medics report that evacuations by day are
rare. Teams wait for darkness, bad weather or smoke, and in some sectors the wounded wait in
forward positions for days \cite{wotr,daylight}. FPV drones and loitering munitions are reported
to cause up to 90 \% of casualties in some sectors \cite{dignitas}. Ground robots and drones are
now used to carry the wounded over the last few kilometres \cite{aspi}.

Three consequences shape our design:
\begin{enumerate}
  \item \textbf{Locating the casualty precisely matters more}, because the evacuation team gets one
        attempt, at night, along a route that must avoid drones.
  \item \textbf{Every transmission is a risk.} Radio emissions are located and targeted, so the
        system must get the most information from the fewest and shortest transmissions.
  \item \textbf{The wounded wait longer}, so time-critical threats (bleeding, airway, hypothermia,
        tourniquet time) need tracking for hours, not minutes.
\end{enumerate}

\nanofig{problem}{The situation Near is built for: GPS denied, drones overhead, a medic who can describe
what he sees but can't give a grid reference.}{Split-panel editorial illustration, 16:9, flat vector style
with subtle texture, plain white background, cyan, amber and red accents. Left panel titled ``GPS DENIED'':
a soldier's wrist GPS showing a glitched screen with ``NO FIX'' and a red jamming wave pattern radiating from
a distant antenna truck. Middle panel titled ``DRONES OVERHEAD'': a small quadcopter silhouette above a
tree line, a dotted red cone below it, two medics crouched under trees with a stretcher. Right panel titled
``WHERE ARE WE?'': the medic holding a radio, a speech bubble with a small castle keep icon, a church icon
and a sun icon with arrows, and a question mark over a blank map. Clean sans-serif labels, exactly these three
titles, no other text, no gore. Colour palette (use only these): navy \#0D1110, cyan \#1AA7C4, amber \#E09A2B, red \#E5484D, green \#3FA877, grey \#8B9A94, plus white for text on dark shapes.}

\section{What already exists and why it isn't enough}

\begin{table}[htbp]\centering\small
\caption{How existing approaches handle a speaker who doesn't know his position.}\label{tab:gap}
\begin{tabularx}{\linewidth}{>{\raggedright}p{3.6cm}X>{\raggedright\arraybackslash}p{3.8cm}}
\toprule
Approach & What it does & With GPS jammed and an unsure speaker\\
\midrule
Device GPS in tactical apps (ATAK-type tools, battle management systems) & Shares the position of each device & No fix, or a spoofed one\\
Other teams on this challenge (e.g. LowKey, Ground Truth, ELARIS) & Speech $\rightarrow$ place name $\rightarrow$ geocoded pin; trust scores and ageing of reports & Needs the speaker to name the place he is at\\
Voice-guided navigation (hasky labs) & Gives directions from a known position & Needs the starting position\\
Acoustic gunshot and drone sensors & Locate sources with dedicated microphone arrays & New hardware, fixed installations\\
Map reading and compass resection & The trained soldier's manual method & Slow, needs a map and calm; few do it under fire\\
\bottomrule
\end{tabularx}
\end{table}

The manual method in the last row is the one we automate. Resection from bearings to known points
is standard land-navigation doctrine, and Pothenot's problem is its classic form. Near does it
from speech, tolerates noise and missing data, and asks for the observation that would help most.

\section{Requirements we set ourselves}

\begin{itemize}
  \item No GPS, no internet, no phone network required. Any of them is used if present.
  \item No new hardware in the field. The soldier's radio is enough; a phone adds sensors.
  \item Answers in seconds: under 300 ms per position update, a few seconds per report.
  \item Honest uncertainty: every position comes with a radius and, when several areas fit, all of them.
  \item English and French, and radio procedure words (callsigns, ``over'').
  \item Medical output is decision support for trained medics: organising, timing and ordering, never diagnosing.
\end{itemize}
